\documentclass[a4paper,12pt]{article}
\usepackage[italian]{babel}
\usepackage[utf8]{inputenc}
\usepackage[T1]{fontenc}
\usepackage{geometry}
\geometry{margin=2.5cm}

\title{Verifica di Basi di Dati -- SQL}
\author{Prof. Fedeli Massimo -- Fabbrica Digitale}
\date{16/03/2026}

\begin{document}
	
	\maketitle
	
	\section*{Schema Logico}
	
	PAZIENTI(\underline{id\_paziente}, nome, cognome, data\_nascita, citt\`a)
	
	MEDICI(\underline{id\_medico}, nome, cognome, specializzazione)
	
	REPARTI(\underline{id\_reparto}, nome\_reparto, piano, id\_medico)
	
	RICOVERI(\underline{id\_paziente}, \underline{id\_reparto}, data\_ricovero)
	
	VISITE(\underline{id\_paziente}, \underline{id\_medico}, voto, data\_visita)
	
	
	\section*{Parte 1 -- Query SQL}
	
	Scrivere le query SQL per le seguenti richieste.
	
	\begin{enumerate}
		
		\item Elencare nome e cognome di tutti i pazienti.
		
		\item Elencare tutti i reparti con il relativo piano.
		
		\item Elencare i pazienti residenti a Roma.
		
		\item Elencare i pazienti ordinati per cognome e nome.
		
		\item Elencare i reparti con il nome del medico responsabile.
		
		\item Elencare i pazienti che hanno effettuato almeno una visita.
		
		\item Elencare i pazienti che hanno ricevuto almeno una valutazione maggiore o uguale a 8.
		
		\item Calcolare la valutazione media per ogni reparto.
		
		\item Elencare i pazienti che hanno effettuato pi\`u di 2 visite.
		
		\item Elencare i medici che seguono pi\`u di un reparto.
		
		\item Elencare i reparti che non hanno ancora visite registrate.
		
		\item Trovare il paziente con la media delle valutazioni pi\`u alta.
		
		\item Elencare i pazienti ricoverati in tutti i reparti a cui sono assegnati.
		
		\item Elencare il numero di pazienti ricoverati in ciascun reparto.
		
		\item Trovare il reparto con la valutazione media pi\`u alta.
		
	\end{enumerate}
	
	
	\section*{Parte 2 -- INSERT}
	
	Inserire un nuovo paziente:
	Mario Rossi, nato il 15 marzo 2000, residente a Bologna.
	
	
	\section*{Parte 3 -- UPDATE}
	
	Aggiornare il piano del reparto ``Cardiologia'' portandolo al piano 3.
	

	\section*{Parte 4 -- SQL DDL}
	
	Scrivere l'istruzione SQL per creare la tabella:
	
	SALE(\underline{id\_sala}, edificio, capienza)
	
	Vincoli:
	\begin{itemize}
		\item \texttt{id\_sala} chiave primaria
		\item \texttt{capienza} $> 0$
	\end{itemize}
		
	\section*{Parte 5 -- Viste}
	
	Creare una vista chiamata \texttt{media\_pazienti} che mostri:
	\begin{itemize}
		\item \texttt{id\_paziente}
		\item \texttt{nome}
		\item \texttt{cognome}
		\item media delle valutazioni
	\end{itemize}
	
	
	\newpage
	\section*{SOLUZIONI}
	
	\subsection*{Parte 1 -- Query SQL}
	
	\begin{enumerate}
		
\item Elencare nome, cognome e città dei pazienti nati dopo
il 1980, ordinati per data di nascita decrescente.
\begin{verbatim}
	SELECT nome, cognome, citta
	FROM PAZIENTI
	WHERE data_nascita > '1980-12-31'
	ORDER BY data_nascita DESC;
\end{verbatim}

\item Elencare i reparti ai piani 2 o 3 con il numero
di pazienti ricoverati.
\begin{verbatim}
	SELECT R.nome_reparto, R.piano,
	COUNT(RC.id_paziente) AS num_pazienti
	FROM REPARTI R
	LEFT JOIN RICOVERI RC ON R.id_reparto = RC.id_reparto
	WHERE R.piano IN (2, 3)
	GROUP BY R.id_reparto, R.nome_reparto, R.piano;
\end{verbatim}

\item Elencare i pazienti di Roma o Milano con almeno
una visita con valutazione superiore a 7.
\begin{verbatim}
	SELECT DISTINCT P.nome, P.cognome
	FROM PAZIENTI P
	JOIN VISITE V ON P.id_paziente = V.id_paziente
	WHERE P.citta IN ('Roma', 'Milano')
	AND V.voto > 7;
\end{verbatim}
		
		\item Elencare i pazienti ordinati per cognome e nome.
		\begin{verbatim}
			SELECT nome, cognome
			FROM PAZIENTI
			ORDER BY cognome, nome;
		\end{verbatim}
		
		\item Elencare i reparti con il nome del medico responsabile.
		\begin{verbatim}
			SELECT R.nome_reparto, M.nome, M.cognome
			FROM REPARTI R
			JOIN MEDICI M ON R.id_medico = M.id_medico;
		\end{verbatim}
		
		\item Elencare i pazienti che hanno effettuato almeno una visita.
		\begin{verbatim}
			SELECT DISTINCT P.nome, P.cognome
			FROM PAZIENTI P
			JOIN VISITE V ON P.id_paziente = V.id_paziente;
		\end{verbatim}
		
		\item Elencare i pazienti che hanno ricevuto almeno una
		valutazione maggiore o uguale a 8.
		\begin{verbatim}
			SELECT DISTINCT P.nome, P.cognome
			FROM PAZIENTI P
			JOIN VISITE V ON P.id_paziente = V.id_paziente
			WHERE V.voto >= 8;
		\end{verbatim}
		
		\item Calcolare la valutazione media per ogni reparto.
		\begin{verbatim}
			SELECT R.nome_reparto, AVG(V.voto) AS media_valutazione
			FROM REPARTI R
			JOIN VISITE V ON R.id_reparto = V.id_medico
			GROUP BY R.id_reparto, R.nome_reparto;
		\end{verbatim}
		
		\item Elencare i pazienti che hanno effettuato più di 2 visite.
		\begin{verbatim}
			SELECT P.nome, P.cognome, COUNT(*) AS num_visite
			FROM PAZIENTI P
			JOIN VISITE V ON P.id_paziente = V.id_paziente
			GROUP BY P.id_paziente, P.nome, P.cognome
			HAVING COUNT(*) > 2;
		\end{verbatim}
		
		\item Elencare i medici che seguono più di un reparto.
		\begin{verbatim}
			SELECT M.nome, M.cognome, COUNT(*) AS num_reparti
			FROM MEDICI M
			JOIN REPARTI R ON M.id_medico = R.id_medico
			GROUP BY M.id_medico, M.nome, M.cognome
			HAVING COUNT(*) > 1;
		\end{verbatim}
		
		\item Elencare i reparti che non hanno ancora visite registrate.
		\begin{verbatim}
			SELECT R.nome_reparto
			FROM REPARTI R
			LEFT JOIN RICOVERI RC ON R.id_reparto = RC.id_reparto
			LEFT JOIN VISITE V ON RC.id_paziente = V.id_paziente
			WHERE V.id_paziente IS NULL;
		\end{verbatim}
		
		\item Trovare il paziente con la media valutazioni più alta.
		\begin{verbatim}
			SELECT P.nome, P.cognome, AVG(V.voto) AS media
			FROM PAZIENTI P
			JOIN VISITE V ON P.id_paziente = V.id_paziente
			GROUP BY P.id_paziente, P.nome, P.cognome
			ORDER BY media DESC
			LIMIT 1;
		\end{verbatim}
		
		\item Elencare i pazienti ricoverati in tutti i reparti
		a cui sono assegnati.
		\begin{verbatim}
					SELECT P.nome, P.cognome
					FROM PAZIENTI P
					JOIN RICOVERI RC ON P.id_paziente = RC.id_paziente
					JOIN VISITE V ON P.id_paziente = V.id_paziente
					GROUP BY P.id_paziente, P.nome, P.cognome
					HAVING COUNT(DISTINCT V.id_medico) = (
					SELECT COUNT(*)
					FROM RICOVERI RC2
					WHERE RC2.id_paziente = P.id_paziente
					);
		\end{verbatim}
		
		\item Elencare il numero di pazienti ricoverati
		in ciascun reparto.
		\begin{verbatim}
			SELECT R.nome_reparto, COUNT(RC.id_paziente) AS num_pazienti
			FROM REPARTI R
			LEFT JOIN RICOVERI RC ON R.id_reparto = RC.id_reparto
			GROUP BY R.id_reparto, R.nome_reparto;
		\end{verbatim}
		
		\item Trovare il reparto con la valutazione media più alta.
		\begin{verbatim}
			SELECT R.nome_reparto, AVG(V.voto) AS media
			FROM REPARTI R
			JOIN RICOVERI RC ON R.id_reparto = RC.id_reparto
			JOIN VISITE V ON RC.id_paziente = V.id_paziente
			GROUP BY R.id_reparto, R.nome_reparto
			ORDER BY media DESC
			LIMIT 1;
		\end{verbatim}
		
	\end{enumerate}
	
	\noindent\rule{\linewidth}{0.4pt}
	
	\subsection*{Parte 2 -- INSERT}
	
	\begin{verbatim}
		INSERT INTO PAZIENTI (nome, cognome, data_nascita, città)
		VALUES ('Mario', 'Rossi', '2000-03-15', 'Bologna');
	\end{verbatim}
	
	\noindent\rule{\linewidth}{0.4pt}
	
	\subsection*{Parte 3 -- UPDATE}
	
	\begin{verbatim}
		UPDATE REPARTI
		SET piano = 3
		WHERE nome_reparto = 'Cardiologia';
	\end{verbatim}
	
	\noindent\rule{\linewidth}{0.4pt}
	
	\subsection*{Parte 4 -- SQL DDL}
	
\begin{verbatim}
	CREATE TABLE SALE (
	id_sala     INT          PRIMARY KEY,
	edificio    VARCHAR(50)  NOT NULL,
	capienza    INT          NOT NULL CHECK (capienza > 0)
	);
	
	CREATE TABLE INTERVENTI (
	id_intervento    INT          PRIMARY KEY,
	descrizione      VARCHAR(100) NOT NULL,
	data_intervento  DATE         NOT NULL,
	id_sala          INT          NOT NULL,
	id_medico        INT          NOT NULL,
	FOREIGN KEY (id_sala)    REFERENCES SALE(id_sala),
	FOREIGN KEY (id_medico)  REFERENCES MEDICI(id_medico)
	);
\end{verbatim}
	
	\noindent\rule{\linewidth}{0.4pt}
	
	\subsection*{Parte 5 -- Viste}
	
	\begin{verbatim}
		CREATE VIEW media_pazienti AS
		SELECT P.id_paziente, P.nome, P.cognome,
		AVG(V.voto) AS media_valutazioni
		FROM PAZIENTI P
		JOIN VISITE V ON P.id_paziente = V.id_paziente
		GROUP BY P.id_paziente, P.nome, P.cognome;
	\end{verbatim}
	
	
	
\end{document}